\documentclass{article}
\usepackage[utf8]{inputenc}
\usepackage[spanish]{babel}
\usepackage{amsmath}
\usepackage{xcolor}
\usepackage{hyperref}

% Definición de estilos para resaltar palabras clave
\newcommand{\keyword}[1]{\colorbox{gray!15}{\textcolor{red!70!black}{\texttt{#1}}}}
\newcommand{\codevar}[1]{\colorbox{gray!15}{\texttt{#1}}}
\newcommand{\inlinebox}[1]{\colorbox{gray!15}{\texttt{#1}}}

\begin{document}

\section*{Pregunta técnica Python}

\subsection*{Definiciones}

Un \keyword{portafolio} de inversión está compuesto por una combinación de $N$ \keyword{activos}. En cada instante de tiempo $t$ el monto en dólares para el activo $i$ viene dado por la variable $x_{i,t}$ y por tanto el valor total del portafolio en el tiempo $t$ ($V_t$) es equivalente a:

\[ V_t = \sum_{i=1}^{N} x_{i,t} \]

El \keyword{precio} de cada activo en el tiempo viene dado por $p_{i,t}$ y la \keyword{cantidad} de cada activo en el tiempo viene dada por $c_{i,t}$, de forma tal que:

\[ x_{i,t} = p_{i,t} \cdot c_{i,t} \]

A su vez se conoce como \keyword{weight} ($w_{i,t}$) de cada activo al \% que representa este sobre el portafolio total. En otras palabras, se cumple que:

\begin{equation}
    w_{i,t} = \frac{x_{i,t}}{V_t} = \frac{p_{i,t} \cdot c_{i,t}}{V_t}
\end{equation}

Considere un portafolio que tiene su inicio en $t = 0$ con valor inicial $V_0$. La cantidad invertida por activo ($c_{i,0}$) se puede calcular usando las definiciones anteriores lo que da como resultado:

\begin{equation}
    c_{i,0} = \frac{w_{i,0} \cdot V_0}{p_{i,0}}
\end{equation}

En el archivo \textbf{datos.xlsx} se encuentran dos sets de datos separados por hoja:
\begin{itemize}
    \item \textbf{Weights:} valores para $w_{i,0}$ de dos portafolios (1 y 2 en las columnas C y D respectivamente) donde el tiempo $t = 0$ equivale al 15/2/22 e $i = 1, \dots, 17$ corresponde a cada uno de los 17 activos invertibles.
    \item \textbf{Precios:} valores para $p_{i,t}$ donde cada columna corresponde a cada uno de los 17 activos y cada fila corresponde al tiempo $t = \text{15/2/22}, \dots, \text{16/2/23}$.
\end{itemize}

\section*{Preguntas}

\begin{enumerate}
    \item Cree un proyecto en \keyword{Django} que permita modelar la definición anterior. Esto es, activos, portafolios, precios, cantidades, weights, montos y cualquier otro elemento que estime conveniente. Considere que el proyecto debe permitir responder las siguientes preguntas.
    \item Genere una función tipo ETL que permita leer y cargar los datos del archivo \textbf{datos.xlsx} a la base de datos correspondiente al proyecto Django mencionado en el punto anterior.
    \item Considere que tanto el portafolio 1 como el portafolio 2 tienen un valor inicial al 15/02/22 ($V_0$) de \$1,000,000,000. Calcule las cantidades iniciales ($c_{i,0}$) para cada uno de los 17 activos en cada uno de los 2 portafolios.
    \item A partir del 15/02/22 los valores de las cantidades se mantienen invariantes, $c_{i,t} = c_{i,0}$, y por tanto los valores de $x_{i,t}$, $w_{i,t}$ y $V_t$ evolucionan debido a la variación en el tiempo de los precios $p_{i,t}$ y las definiciones explicitadas en el inicio del documento. Genere endpoints tipo API REST que reciban los parámetros \codevar{fecha\_inicio} y \codevar{fecha\_fin} y entregue los valores entre esas fechas para $w_{i,t}$ y $V_t$. Se espera el uso del ORM de Django para obtener los datos necesarios para los cálculos.
    \item Genere un \codevar{view} que utilice la API anterior donde se pueda comparar de manera gráfica la evolución en el tiempo de las variables $w_{i,t}$ y $V_t$. Para $w_{i,t}$ se recomienda un gráfico tipo ``stacked area'' y para $V_t$ gráficos de línea.
    \item \textbf{Bonus}: Estructurar el proyecto de Django siguiendo la siguiente guía de estilos: \url{https://github.com/HackSoftware/Django-Styleguide}.
\end{enumerate}

\end{document}